% ============================================================================
%  SECCIÓN 3.16: README DEL PROYECTO
% ============================================================================

\section{README del proyecto}

El archivo README es el documento de presentaci\'on de un repositorio de
software, en el que se describe el prop\'osito del proyecto, los requisitos de
instalaci\'on y la forma de ejecuci\'on, de modo que cualquier persona pueda
comprender y poner en marcha la aplicaci\'on \citep{githubdocs}.

El README del proyecto se elabor\'o en el lenguaje Markdown e incluye los
siguientes apartados:

\begin{itemize}
  \item T\'itulo del proyecto y breve descripci\'on de la aplicaci\'on
    \nombreApp{} y su prop\'osito.

  \item Requisitos previos: sistema operativo compatible, Node.js y la
    aplicaci\'on Expo Go en el dispositivo m\'ovil.

  \item Instrucciones de instalaci\'on, que describen la descarga del
    repositorio y la instalaci\'on de las dependencias del proyecto.

  \item Instrucciones de ejecuci\'on, que indican c\'omo iniciar el servidor de
    desarrollo de Expo y c\'omo abrir la aplicaci\'on en el dispositivo.

  \item Enlace al repositorio del proyecto en GitHub
    (\url{https://github.com/Joaquinfr87/proyecto-movil}).
\end{itemize}

La siguiente figura muestra la estructura general del contenido del archivo
README del proyecto.

\begin{figure}[H]
\centering
  \figuraAPA{fig:readme-proyecto}{Contenido del archivo README del proyecto}{%
    \includegraphics[width=0.62\linewidth]{readme.jpeg}%
  }{Elaboraci\'on propia (captura del archivo README).}
\end{figure}

\subsection{Secci\'on de persistencia de datos en el README}

Como parte del avance~5, el README del repositorio se actualiz\'o con una
secci\'on titulada \textit{Persistencia de datos}. Esta secci\'on
explica:

\begin{itemize}
  \item \textbf{Qu\'e informaci\'on se almacena:} la sesi\'on del usuario
    (token JWT y refresh token) de forma local mediante AsyncStorage, y
    los datos de dominio (escenarios, eventos, favoritos y perfil) en
    Supabase.

  \item \textbf{Por qu\'e se almacena:} para conservar la sesi\'on y los
    favoritos entre aperturas de la aplicaci\'on, evitando que el usuario
    tenga que iniciar sesi\'on o guardar sus preferencias
    nuevamente.

  \item \textbf{Qu\'e tecnolog\'ia se utiliz\'o:} AsyncStorage para el
    almacenamiento local y Supabase (PostgreSQL con pol\'iticas RLS)
    para la base de datos remota.

  \item \textbf{Qu\'e ocurre al cerrar la aplicaci\'on:} la sesi\'on
    persiste localmente y los datos de dominio permanecen en el servidor,
    por lo que pueden recuperarse al abrir la aplicaci\'on nuevamente.

  \item \textbf{Qu\'e informaci\'on puede eliminarse:} los favoritos,
    los escenarios (soft/hard delete) y los eventos, que el usuario o
    el administrador puede gestionar desde la aplicaci\'on.
\end{itemize}

Este apartado refleja la evoluci\'on del proyecto hacia una capa de
persistencia documentada y reproducible por cualquier colaborador del
repositorio.
